\subsection{Connected sums} \label{sec-connected-sum}

Consider the decomposition of an $(n-1)$-dimensional pseudo-manifold $\Delta$ as the connected sum of two pseudo-manifolds
\[ \Delta = \Delta_1 \#_D \Delta_2.\]
Here $D$ is a common $(n-1)$-dimensional subcomplex of $\Delta_1$ and  $\Delta_2$ that is a pseudo-manifold with boundary. We remove the interior of $D$ from $\Delta_1, \Delta_2$, and glue the remaining complexes along their common boundary. We assume that $\Delta$, $\Delta_1$ and $\Delta_2$ are oriented compatibly. This means that if a maximal simplex lies in both $\Delta$ and $\Delta_i$, then it has the same orientation in both.


  \begin{figure}[htb]
    \centering
    \includegraphics[width=0.8\textwidth]{figures/connected_sum.pdf}
    \caption{Connected sum of $1$-dimensional spheres $\Delta_1$ and $\Delta_2$ along $D$.}
  \end{figure}
  
  Let us also denote the simplicial complexes $\Delta_1$ and  $\Delta_2$ glued along $D$ by 
\[ \tilde{\Delta} = \Delta_1 \cup_D \Delta_2.\]
Let $\tilde{\theta}_i \in \cA^1(\tilde\Delta)$ be linear parameters for $\tilde{\Delta}$. These parameters, viewed as piecewise linear functions on $\tilde{\Delta}$, restrict to linear parameters on $\Delta$, $\Delta_1$ and $\Delta_2$. Similarly, a piecewise polynomial function $\tilde{f} \in \cA^n(\tilde{\Delta})$ restricts to piecewise polynomial functions $\tilde{f}|_\Delta \in \cA^n(\Delta)$,  $\tilde{f}|_{\Delta_1} \in \cA^n(\Delta_1)$ and $\tilde{f}|_{\Delta_2} \in \cA^n(\Delta_2)$. The latter two agree on $D$.

\begin{lemma} 
Let $\tilde{f} \in \cA^n(\tilde{\Delta})$. Then
\[ \pi_{\Delta} (\tilde{f}|_\Delta) = \pi_{\Delta_1}(\tilde{f}|_{\Delta_1}) + \pi_{\Delta_2}(\tilde{f}|_{\Delta_2}).\]
\end{lemma} 

\begin{proof} 
The maximal simplices of $D$ appear in $\Delta_1$ and $\Delta_2$ with opposite orientations. Hence these terms cancel on the right hand side. The remaining terms give the left hand side.
\end{proof}

\begin{rema} 
The  previous lemma was used in \cite{Karu, BL2, BBFK2}. Its meaning as integration over a connected sum was realized by Karl-Heinz Fieseler. The lemma says that Brion's integration map behaves like ordinary integration. One can decompose the domain of integration into pieces and sum the integrals over the pieces.
\end{rema}

We next consider a more general connected sum. Let $v_0$ be a new vertex and let 
\[ C(\Delta) = \{v_0\} * \Delta \]
be the cone over $\Delta$ with vertex $v_0$.  Let $\pi_i = \{v_0\} * \sigma_i$, $i=1,\ldots,M$ be the maximal simplices in $C(\Delta)$, and let $\Pi_i = \partial \pi_i$ be the simplicial $(n-1)$-spheres. Then
\[ \Delta = \#_{i=1}^M \Pi_i.\]
Here we use a more general notion of connected sum. We assume that $\Delta$ and $\Pi_i$ are oriented compatibly. Then an $(n-1)$-simplex $\sigma \in \Delta$ appears in the disjoint union $\sqcup_{i=1}^M \Pi_i$ exactly once and with the same orientation as in $\Delta$. All other $(n-1)$-simplices of  $\sqcup_{i=1}^M \Pi_i$ appear there twice with opposite orientations.

  \begin{figure}[htb]
    \centering
    \def\svgwidth{0.8\columnwidth}
    \includegraphics[width=0.8\textwidth]{figures/general_connected_sum.pdf}
    \caption{Decomposition of a $1$-sphere as a connected sum.}
  \end{figure}

As before, we let $\tilde{\Delta}$ be the union of $\Pi_i$. A system of linear parameters $\tilde{\theta}_i$ on $\tilde{\Delta}$ restricts to a system of parameters on $\Delta$ and all $\Pi_i$. 

\begin{lemma} \label{lem-int-sum}
Let $\tilde{f}\in \cA^n(\tilde{\Delta})$. Then 
\[ \pi_{\Delta}(\tilde{f}|_\Delta)  = \sum_{i=1}^M \pi_{\Pi_i}(\tilde{f}|_{\Pi_i}).  \]
\end{lemma}

The parameters $\tilde{\theta}_i$  have extra variables $a_{i,0}$ corresponding to the new vertex $v_0$. We may include these in the field $K$,
\[ K = k(\a{i, j})_{i=1,\ldots,n; j=0,\ldots, N}.\]  
However, the  map $\pi_{\Delta}$ does not depend on the variables $\a{i, 0}$. If $f\in \cA^n(\Delta)$ has coefficients in $k(\a{i, j})_{i=1,\ldots,n; j=1,\ldots, N}$ then $\pi_{\Delta}(f)$ also lies in the same field.


An alternative connected sum decomposition would be to take one of the existing vertices, say $v_1$, as the cone point and replace $C(\Delta)$ with 
\[ \{v_1\}* (\Delta \setminus \Star^\circ v_1).\]

  \begin{figure}[htb]
    \centering
  \def\svgwidth{0.8\columnwidth}
    \includegraphics[width=0.8\textwidth]{figures/alt_connected_sum.pdf}
    \caption{Alternative decomposition of a $1$-sphere as a connected sum.}
  \end{figure}
